\section{Reset de equipos de red}

Antes de proceder con la configuración de la infraestructura de red diseñada, fue necesario realizar el proceso de restauración de los dispositivos de red a su estado inicial de fábrica. Este procedimiento garantiza un entorno limpio, libre de configuraciones previas que puedan interferir con el correcto funcionamiento de la red o afectar los resultados de las pruebas de seguridad realizadas.

El reseteo de los equipos permite asegurar que todas las configuraciones aplicadas durante la práctica sean controladas, documentadas y reproducibles, lo cual constituye una buena práctica fundamental en la administración de redes y en auditorías de seguridad informática.

\subsection{Reset del Router}

El router Cisco fue restaurado a su configuración inicial con el objetivo de eliminar cualquier parámetro previamente almacenado, incluyendo configuraciones de interfaces, rutas, contraseñas y servicios. Para este propósito, se utilizó el procedimiento estándar que permite ignorar el archivo de configuración almacenado en la memoria NVRAM durante el arranque del dispositivo.

Este proceso se realiza mediante la modificación temporal del registro de configuración del router, permitiendo acceder al sistema sin cargar la configuración existente.

\subsubsection{Procedimiento de reseteo}

A continuación, se detallan los comandos ejecutados durante el proceso de reseteo del router:

\begin{lstlisting}[caption={Comandos utilizados para el reset del Router Cisco}]
confreg 0x2142
reset
no
enable
config t
config-register 0x2102
end
copy running-config startup-config
reload
no
\end{lstlisting}

\subsubsection{Descripción del procedimiento}

El comando \texttt{confreg 0x2142} modifica el registro de configuración del router para que ignore el archivo \texttt{startup-config} durante el arranque. Posteriormente, el comando \texttt{reset} reinicia el dispositivo, permitiendo el acceso al modo privilegiado sin aplicar configuraciones previas.

Una vez iniciado el router, se accede al modo de configuración global para restaurar el registro de configuración original mediante el comando \texttt{config-register 0x2102}, asegurando que en futuros reinicios el router cargue correctamente la configuración almacenada. Finalmente, se guarda la configuración limpia en la memoria NVRAM y se reinicia el dispositivo para aplicar los cambios de forma permanente.

\subsection{Reset del Switch}

El reseteo del switch Cisco se realizó con el objetivo de eliminar configuraciones anteriores relacionadas con VLANs, puertos, contraseñas y otros parámetros administrativos. Este procedimiento es especialmente importante en dispositivos de capa 2, ya que las VLANs se almacenan en un archivo independiente (\texttt{vlan.dat}), el cual debe ser eliminado manualmente.

\subsubsection{Procedimiento de reseteo}

A continuación, se presentan los comandos utilizados para restaurar el switch a su estado inicial:

\begin{lstlisting}[caption={Comandos utilizados para el reset del Switch Cisco}]
enable
show flash
delete vlan.dat
--confirm
erase startup-config
reload
--no
--no
\end{lstlisting}

\subsubsection{Descripción del procedimiento}

El proceso inicia accediendo al modo privilegiado del switch mediante el comando \texttt{enable}. A continuación, se verifica el contenido de la memoria flash para confirmar la existencia del archivo \texttt{vlan.dat}, el cual almacena la información de las VLANs configuradas previamente.

El comando \texttt{delete vlan.dat} elimina dicho archivo, asegurando que no existan VLANs residuales que puedan afectar la segmentación de la red. Posteriormente, el comando \texttt{erase startup-config} borra la configuración de inicio almacenada en la NVRAM. Finalmente, el switch es reiniciado mediante el comando \texttt{reload}, descartando cualquier solicitud de guardar configuraciones, lo que permite iniciar el dispositivo con una configuración completamente limpia.

\subsection{Importancia del reseteo en el proyecto}

El reseteo de los dispositivos de red constituye una etapa crítica dentro del desarrollo del proyecto, ya que garantiza que todas las configuraciones aplicadas posteriormente respondan exclusivamente a los requerimientos del diseño propuesto por el Grupo 6. Además, este procedimiento facilita la identificación de errores, mejora la trazabilidad de los cambios realizados y asegura que los resultados obtenidos durante las pruebas de seguridad sean confiables y reproducibles.

La correcta ejecución de estos procesos sienta las bases para una configuración ordenada de la red y para la realización de las actividades de monitoreo y ataque documentadas en los capítulos siguientes.
